% Generated by report/tools/tables.py from the result files of the repository. Do not edit.
\begin{table}[tbp]
\centering
\footnotesize
\setlength{\tabcolsep}{2pt}
\caption[Grille v2 : 2018 contre RL classique 2026]{Grille v2 (critères Y1 à Y12, écarts 11 et 14 de \file{study/PROTOCOL.md}, écrite avant toute lecture v2), sur les mêmes échantillons que le tableau~\ref{tab:sondage-pineau} : 50 articles de 2018 lus par nous et 50 articles de RL classique de 2026. Mêmes colonnes : $k/n$ sur les articles où le critère s'applique, part avec intervalle de Wilson à 95~\%, différence en points avec l'intervalle de Newcombe à 95~\%, test exact de Fisher, correction de Holm sur les douze critères de la grille v2 (famille distincte de la grille v1). Y3 vaut « oui » quand la conclusion principale porte sur au moins deux tâches ; Y2 et Y11 ne s'appliquent qu'aux articles qui rapportent une incertitude, Y4 qu'aux articles à plusieurs tâches. \textbf{En gras}, les critères dont la différence passe la correction de Holm : seul Y9. Valeurs harmonisées (\file{study/overrides\_v2.json}). Source : \file{study/extension\_stats.json}.}
\label{tab:sondage-extension}
\begin{tabularx}{\linewidth}{@{}l>{\raggedright\arraybackslash}Xr r@{\hspace{0.28em}}l r r@{\hspace{0.28em}}l r@{\hspace{0.28em}}l rr@{}}
\toprule
 & & \multicolumn{3}{c}{\makecell{Notre lecture\\{}de 2018}} & \multicolumn{3}{c}{\makecell{RL classique\\{}2026}} & \multicolumn{2}{c}{Différence} & &  \\
\cmidrule(lr){3-5}\cmidrule(lr){6-8}
Code & Critère & $k/n$ & \multicolumn{2}{c}{\% [Wilson]} & $k/n$ & \multicolumn{2}{c}{\% [Wilson]} & \multicolumn{2}{c}{[Newcombe]} & $p$ & $p$ Holm \\
\midrule
Y1 & Mesure de performance définie & 38/48 & 79 & [66\,;\,88] & 39/50 & 78 & [65\,;\,87] & −1 & [−17\,;\,+15] & 1 & 1 \\
Y2 & Barres d'erreur définies & 25/30 & 83 & [66\,;\,93] & 38/44 & 86 & [73\,;\,94] & +3 & [−13\,;\,+21] & 0,75 & 1 \\
Y3 & Plusieurs tâches (au moins 2) & 43/48 & 90 & [78\,;\,95] & 41/50 & 82 & [69\,;\,90] & −8 & [−22\,;\,+7] & 0,39 & 1 \\
Y4 & Une tâche perdue montrée & 34/43 & 79 & [65\,;\,89] & 38/41 & 93 & [81\,;\,97] & +14 & [−2\,;\,+29] & 0,12 & 1 \\
Y5 & Origine du code des références & 22/48 & 46 & [33\,;\,60] & 20/49 & 41 & [28\,;\,55] & −5 & [−24\,;\,+14] & 0,68 & 1 \\
Y6 & Exclusions dites & 14/48 & 29 & [18\,;\,43] & 16/50 & 32 & [21\,;\,46] & +3 & [−15\,;\,+20] & 0,83 & 1 \\
Y7 & Version des bancs d'essai & 11/48 & 23 & [13\,;\,37] & 13/50 & 26 & [16\,;\,40] & +3 & [−14\,;\,+20] & 0,82 & 1 \\
Y8 & Test de généralisation & 11/48 & 23 & [13\,;\,37] & 17/50 & 34 & [22\,;\,48] & +11 & [−7\,;\,+28] & 0,27 & 1 \\
\textbf{Y9} & \textbf{Sensibilité à un hyperparamètre} & \textbf{21/48} & \textbf{44} & \textbf{[31\,;\,58]} & \textbf{40/50} & \textbf{80} & \textbf{[67\,;\,89]} & \textbf{+36} & \textbf{[+17\,;\,+52]} & \textbf{0,0003} & \textbf{0,004} \\
Y10 & Versions des bibliothèques & 1/48 & 2 & [0\,;\,11] & 4/50 & 8 & [3\,;\,19] & +6 & [−4\,;\,+17] & 0,36 & 1 \\
Y11 & Type d'incertitude nommé & 25/30 & 83 & [66\,;\,93] & 41/44 & 93 & [82\,;\,98] & +10 & [−5\,;\,+27] & 0,26 & 1 \\
Y12 & Outil statistique postérieur à 2018 & 1/48 & 2 & [0\,;\,11] & 7/50 & 14 & [7\,;\,26] & +12 & [+1\,;\,+24] & 0,06 & 0,66 \\
\bottomrule
\end{tabularx}
\end{table}
